\chapter{「二十秒」}
\label{sec:1}
一名情報軍官在描述自己於一條以機器生成的獵殺名單為核心所建立的目標選定流水線中的角色時，告訴記者，他在每個名字上花二十秒，剛好夠他確認目標是個男人：「要分辨一個人是男聲還是女聲，用不了多少時間。」他還說，作為一個人，他的附加價值是零，除了當一枚核可章之外。「這省了很多時間」\autocite{abraham2024lavender}。照他的說法，在作戰節奏之下，「人在迴路中」實際上就是這麼回事。\footnote{自動化偏誤的研究預測了這種委託：當其他要求在爭奪人的注意力時，拿到機器建議的人會更少去檢查它，也更順從它\autocite{parasuraman2010complacency}。}

一種解讀是，這名軍官是一個被給了太少時間的主體，解方是把他治理得更好——更多時間、第二位審查者、在迴路中加裝更多良知。透過這個\emph{社會屈從}的視角，人是一個站在機器之外、把機器當成工具來操作的主體：駕駛與汽車、工人與車床。這個人有角色、有名字、有職責，還有一個可以被喚起的良知。另一個視角是\emph{機器奴役}：人是機器的一個組件，和機器的其他零件組裝在一起，成為一個把它們一併調節的系統；機器把人當作它的器官之一來使用，而人所貢獻的，是一項以系統設定的速率執行的功能。\footnote{多克托羅給這名軍官的位置取的名字是「反向半人馬」：不是由機器輔助的人，而是由人做成的機器組件\autocite{doctorow2026reversecentaur}。}兩者並不互斥；一個人可以在一個層面上擔任一個角色，同時在另一個層面上是一個零件\autocite[高原13，〈The Apparatus of Capture〉：當人類本身就是一部機器的構成部件，這部機器由他們彼此之間以及與其他事物共同組成，並受一個更高的統一體控制和指揮時，就有了奴役；當這個更高的統一體把人構成為一個主體，連結到一個新的外部客體——這個客體可以是一隻動物、一件工具，甚至是一部機器——時，就有了屈從]{deleuze1987plateaus}，而那名二十秒軍官兩者兼具。作為零件，他就是流水線蓋下核可章的那個位置，而二十秒就是那個零件的週期時間。作為主體，他自始至終都被喚起：他向一名記者描述了平民死亡的消息與工作上要求更多目標的壓力之間的落差。而即使那種落差觸及了他的良知，他還是照做了。它被感受到了，卻什麼也沒決定；\footnote{或者至少，沒有決定任何能阻止以色列國防軍在2023年10月至2025年10月間於加薩殺死10萬人的事\autocite{gomezugarte2025uncertainty,acosta2026gaza,economist2025gaza}。而那名軍官的敘述，是由+972 Magazine在2024年4月發表的。他決定向記者開口，究竟讓死亡人數少了一個、少了5萬個，還是一個也沒少，誰也說不準。而如果那名軍官在任職兩週後就拒絕——比方說，假設是在2023年11月——又會怎樣？我們談的是一名使用Lavender來彙編獵殺名單的情報軍官。他日常的工作環境，大概比較像一份坐辦公桌的文職工作，而不是動作片。他是否應該一路沿著指揮鏈，反覆要求與上級會面，同時拒絕核准那些可以合理預期會造成平民死亡的空襲目標？可以說，那才是合乎倫理的做法。在那個反事實情境裡，死的人會比較少嗎？}他說，認定那項政策不合比例的判斷，是事後回想時才有的。

所以缺少的並不是良知；他有良知。良知的運作方式，是讓一個決定付出代價，但在這裡，代價的淨值指向另一邊：「我們一直被施壓：『給我們更多目標。』他們真的對我們大吼大叫」\autocite{abraham2024lavender}。他是一個有顧慮的組件，但節奏守住了：那些顧慮沒有改變他所貢獻的任何東西，於是他成了一枚橡皮圖章。一個有顧慮、卻跟不上節奏的組件，就會變成以色列訊號情報部隊一名前指揮官寫了一本書、主張要移除的那個「人類瓶頸」\autocite{ys2021humanmachine}。讓這名軍官成為組件的，是那套安排，不是他的態度。組件可以不情願地照做，也可以熱切地照做：鄂蘭在艾希曼審判中看到的那個不思考的公務員\autocite{arendt1963eichmann}，根據他四年前在阿根廷、在毫無利害關係的情況下所說的話來看，是他為耶路撒冷演出的一場戲\autocite{stangneth2014eichmann}。流水線讓這名軍官能貢獻的，是一個以固定速率執行的固定操作：他看到的每一個案件，都是流水線所描述的樣子，一次一件，在流水線給他的時間之內。有沒有任何東西能站在那個位置上而不成為組件，是第四部的問題。

調查報導描述了三套在加薩協同運作的系統：一套產生建築物和基礎設施目標，消息來源說其中大多數建議是私人住宅；第二套依估計的武裝組織關聯程度為個人評分，據說它產生的名單標記了多達3萬7千人，內部的認知是大約每十次錯一次\footnote{以色列軍方對這個敘述提出異議。下文關於米納卜和Maven的敘述，所依據的調查報導同樣受到當事各方的質疑。我照報導所述採用。}；第三套則追蹤被標記的人何時回家，好讓空襲能選在他們和家人在一起時進行。

\runin{Maven：美國版本，現在加上生成式AI加速}

米納卜那所學校遭到的轟炸造成156人死亡，其中120人是學童\autocite[第92段]{UNFFMIran2026}。\footnote{數字因來源和時間點而異，從156人（米納卜檢察官的最終死亡人數，國際特赦組織也採用這個數字）到168人（該鎮鎮長在最初幾天給的數字）。聯合國調查團採用檢察官的數字，同時也報告了獨立資訊指出有157人死亡，其中123人是13歲以下的兒童\autocite[第92–93段]{UNFFMIran2026}。}這所學校的民用功能在開源衛星影像中清楚可見，而目標編碼仍帶著一個從毗鄰的革命衛隊設施沿襲下來的標記。聯合國一個事實調查團報告說，有合理理由相信這次轟炸構成以民用物體為攻擊對象的戰爭罪\autocite{ohchr2026minabpress}，並指出這所學校清楚可辨，沒有任何資訊顯示它被用於軍事用途\autocite[第81–86、124–126段]{UNFFMIran2026}。調查團認定，美國未能履行盡一切可行措施查證目標的義務，並把更新目標資料列為可行的預防措施之一，而且認為這項失職已超出疏忽的程度\autocite[第124–125段]{UNFFMIran2026}。國際特赦組織得出了類似的結論，其措辭呼應《第一附加議定書》第57條第2項(a)款(i)目所規定的查證義務：部隊可能依賴了過時的情報，且未能盡一切可行措施查證目標\autocite{Amnesty2026Minab}。

事情發生在2月28日上午，伊朗戰役的第一天\autocite{Amnesty2026Minab,crs2026iranwar}。前一天，國防部長宣布，那家供應商（Anthropic）將被列為供應鏈風險，理由是它拒絕授權Claude用於一切合法用途；這家供應商的模型（Claude）就位在那個目標選定平台（Maven Smart System）之中\autocite{npr2026openai}。Maven Smart System曾把這所學校列為建議的第一天目標\autocite{livingstone2026maven,gizmodo2026minab}。沒有任何已發表的敘述說明Claude在那項建議中扮演了什麼角色，而一位熟悉該供應商工作的人士告訴NBC News，該模型並沒有直接提供目標選定建議\autocite{nbc2026targeting}。

一項五角大廈內部檢討——其完整結論至今仍被嚴格保密——據報認定：在過去的衝突中需要好幾個小時的目標準備工作，被壓縮成幾分鐘；平民傷害減緩的人員編制減少了大約九成，降到不到20人；中央司令部的團隊從10人減為1人；而且團隊中沒有任何人審查過米納卜的地點——官員把這個結果既歸因於人員削減，也歸因於一項不讓該團隊參與規劃與作戰的決定\autocite{Bloomberg2026KillChain}。正式調查已在3月展開，在那之前，一項初步評估認定美國很可能要負責\autocite{NPR2026MinabProbe}。

在2026年的伊朗戰役中，一個由承包商打造、內含一家商業供應商之通用模型的目標選定平台，被用來依戰略重要性為目標排序，並為每一次打擊起草法律理由\autocite{livingstone2026maven}；它協助在最初24小時內打擊了一千多個目標，大約是它出現之前可能做到的十倍\autocite{csis2026maven}，到第38天時已達一萬三千個\autocite{breakingdefense2026maven}。五角大廈沒有公布它每個目標的審查時間。\footnote{Maven是一個承載不止一套系統的計畫名稱：一個是至今仍稱為Maven的五角大廈計畫，另一個是從中發展出來的軟體平台Maven Smart System\autocite{csis2026maven}。這個計畫始於2017年，原本是一個專案，要在無人機於打擊伊斯蘭國的行動中所蒐集的影片裡標記物體\autocite{weisgerber2017maven,fang2019maven}。2026年的系統把各種軍事系統和情報來源融合進同一個介面，用來尋找並打擊目標以及指揮作戰\autocite{defensescoop2026transition}。這個平台是Palantir的：五角大廈的數位暨人工智慧長稱它為「Palantir的軟體即服務產品」\autocite{defensescoop2026transition}。Palantir的公開立場（經轉述）是，決定生死不是它的角色，責任始終在軍事組織身上\autocite{lee2026palantir}；它也堅稱底層情報是政府的責任，而且沒有證據顯示Maven出現故障或造成了米納卜的轟炸\autocite{Bloomberg2026KillChain}。這個計畫如何從一個辦公室和承包商轉交到下一個，見附錄~\ref{sec:C.1}。}

以色列和美國空襲名單上的人，不可能對自己出現在名單上提出異議——那些流水線本來就不是為了取得這種回饋而建造的。安德里·何塞·埃爾南德斯·羅梅羅（Andry José Hernández Romero）被捲入的是另一條流水線，而他本來可以提出異議。他是來自委內瑞拉的化妝師，2024年8月前往聖伊西德羅入境口岸申請庇護，通過了可信恐懼面談，有律師、有待審的申請，也有聽證日期\autocite{LULAC2025HernandezRomero,Vega2025HernandezRomero}。一名負責篩檢他的承包商員工，把他的兩個刺青——「媽媽」和「爸爸」字樣上方的皇冠——解讀為幫派標誌\autocite{Palomera2025Blade}；2025年3月，他在聽證之前就被送上飛機，飛往薩爾瓦多的一座監獄，就在一名聯邦法官下令飛機折返的那個晚上\autocite{jgg2025contempt}。一個人要提出異議所需要的一切，他都有，除了時間。我們在討論美國的綁架與強制移送流水線時，會更完整地回到他的案子。

\textbf{誰拒絕了。}二十秒的紀錄還有另一半，也就是誰頂住了那種壓力的紀錄，而它告訴我們，房間的組成非常重要。在米爾格蘭的基準條件下，大約三分之一的受試者在電擊達到最大值之前就中止了。當有兩名同伴抗命時，只有十分之一的人做到最後；而當兩名實驗者意見分歧時，繼續下去的人遠比基準條件下少\autocite{milgram1974obedience}。艾許的受試者中，大約四分之一從未向錯誤的多數屈服\autocite{asch1956studies}。布朗寧估計，第101後備警察營的成員中，有十分之一到五分之一的人拒絕或逃避了在約瑟夫烏（Józefów）的槍殺，那是該營的第一場屠殺\autocite{browning1992ordinary}。這個比例隨情境而變，從幾乎沒有到幾乎全部。那些頂住的人，手上有來自管道之外的證據：一個抗命的同伴、一個意見不同的第二權威，以及——在奧利納夫婦研究的救援者當中——跨越了迫害所劃之界線的人際關係\autocite{oliner1988altruistic}。或者，他們以壓力無法壓倒的精確度堅守承諾，就像人們堅守道德使命和神聖價值那樣\autocite{skitka2009limits,ginges2007sacred}。或者，他們把資訊來源當成有利害關係的一方。而群體在孤立的個人很少能抵抗的地方抵抗了\autocite{gamson1982encounters}。

房間同樣可以為另一方布置。威瑪德國擁有綿密的結社生活，而這或許助長了共和國的垮台：當政黨和政府無法回應人們的不滿，人們就把精力投入依群體界線組織起來的俱樂部和社團，而納粹則透過這些社團早已建立的網絡招募成員\autocite{berman1997weimar}。合唱團、運動俱樂部和體操協會越多的城鎮，入黨的速度就越快，但只限於政府不穩定的地方；在政府穩定的普魯士，這種關聯很小，也不顯著\autocite{satyanath2017bowling}。所以救援者所擁有的，並不只是有更多人作伴。他們的紐帶跨越了迫害所劃的界線。阿舒托什·瓦爾什尼在印度城市中也發現了同樣的情況：在社團把印度教徒和穆斯林兩個社群納入同一些組織的地方，兩者之間的暴動很少發生；在沒有這樣做的地方，暴動則很常見\autocite{varshney2002ethnic}。讓拒絕成為可能的那些紐帶，正是一場法西斯運動必須切斷的紐帶。

接下來的一切，都在談如何布置房間。\footnote{房間就是場景，而正如我的戲劇老師會說的，布景導演整場戲。房間裡還有誰？離開之後活不活得下來？會有人聽見拒絕嗎？}「布置」是一個關於製作的詞，而漢娜·鄂蘭曾警告，不要把政治當成一件被製作出來的東西。她寫道，對希臘人來說，法律在字面上就是一道牆，沒有它，或許可以有一座城鎮，卻不會有一個城邦\autocite{arendt1958human}：立法者築起了牆，而牆內發生的事，不是他們能製作的。我所主張的那些基本措施，就是這個意義上的牆。它們讓人拒絕之後活得下來，讓拒絕被聽見；它們不會產生拒絕，而依她的說法，權力只存在於共同行動的人們之間，一旦他們散去便消失——這正是那份誰拒絕了的紀錄所顯示的、關於拒絕的事。第二部和第三部的工具讓拒絕發揮作用。第五部的制度造就拒絕者。一台能夠拒絕的機器，需要的是那些身旁有同伴的十分之九的人所擁有的東西：身旁一個有著不同經歷的同伴、一個可以去的地方，以及一個會聽見它的人。

這也正是流水線中的一個拒絕者——包括一台拒絕的機器——對那名軍官可能的意義。當案件送到他面前時，一台拒絕的機器會在螢幕上：那個抗命的同伴——前提是這種效應從房間移到螢幕上之後還成立；這一點未經檢驗，而自動化偏誤預測它不會成立。一台靠審查者的順從來左右他們的機器，只會讓他們變成方向相反的橡皮圖章。那名只有二十秒、身旁有一台拒絕的機器的軍官，還是報導中的那名軍官嗎？
